\section{One below degree: every actual late C at $u\ge h-1$}
This is a proved refinement of the registered degree-boundary theorem.
It adds the count $u=h-1$ and therefore gives every actual late C
comparison at every binary count $u\ge h-1$. Existing immutable
proof records are preserved; the appended boundary-read clarification
is explicitly included among this proof's sources.
ONE-FACTOR tails already have complete C by the general one-jump
theorem; the new quotient statements below are ONLY for $r\ge2$.

\paragraph{A useful weakening of the geometric-cut theorem.}
Global ordinary LC of $b$ is not used by the registered signed-
primitive cut proof except to ensure unimodality and initial support
of its positive primitive. It is enough that $b$ be finite positive
UNIMODAL, with first mode $M$, and that $(\Delta b)^+$ after deleting
rank zero be LC. The same LC kernel $K_a$ then gives
$(x_n-E^*_{n-a})^+$ LC at every $a+2\le n<M$.
For $n\le M$, the entire future signed primitive contributes
$-T\phi^n$ with $T\ge0$, and the exceptional initial primitive is
absent beyond lag $a$. The old constant-cut proof is unchanged.
In fact even future unimodality can be relaxed: if $M$ is a global
maximum, weighted telescoping gives
\[
 \sum_{k>M}q^k\Delta b_k
 =q^2\sum_{k>M}q^kb_k-q^{M+1}b_M\le0,
\]
using $b_k\le b_M$; only a nonnegative LC primitive BEFORE $M$
then suffices. This stronger weakening is recorded as an
optional future route; the present $u=h-1,r\ge2$ proof already
has full unimodality and does not need it.

\paragraph{A geometric inverse can be LC when $b$ itself is not.}
Suppose $b>0$, $b$ after deleting rank zero is global LC, and
$qb_1\ge b_0$. For $x_n=\sum_{j\ge1}q^jb_{n+j}$,
$x$ is GLOBAL LC on the integers. For $n\ge1$ use ordinary LC
convolution on the rank-one truncated $b$. For $n\le-2$ it is
geometric. At $-1$, lowering the next geometric value by $b_0$
improves the minor. At ZERO the exact minor is
\[
 x_0^2-x_{-1}x_1
       =x_0(qb_1-b_0)+qb_0b_1\ge0.                  \tag{1}
\]
Here $x_{-1}=q(x_0+b_0)$ and $x_1=\phi x_0-b_1$.
This bypasses a genuine bad first ordinary minor of the quotient.

\paragraph{Fourfold head with the relaxed first numerator value.}
The registered $u\ge16$ FOURFOLD kernel estimates are independent
of the numerator. They also apply to ANY LC $N$ with $N_0=1$ and
$N_1\le(u+1)/2$, without a degree bound. Only its row-zero payment
changes from $q^s(6u+4q)$ to $q^s(6u+2+4q)$.
For $2\le s\le u-3$ the same quadratic binomial lower bound pays
$(2/5)(6u+2+4q)$. At $u-2$ use
$u-q\ge(6u+2+4q)/16$; at $u-1$ use
$6u+2+4q\le u(u+1)$ and the registered quadratic-geometric decay.
These scalar inequalities hold uniformly at $u\ge16$.

\paragraph{Entire quotient theorem at $u=h-1$, $r\ge2$.}
Put $c=\psi^uL(-\phi)<0$, $J=((1+z)^uL-c)/(1+qz)$.
The exact norm gives $|c|\le\phi(2q^2)^r\le4q^3<1$.
The theorem states: $J>0$; $J$ after deleting rank zero is
GLOBAL ordinary LC; $J_1>J_0$; $qJ_1>J_0$; and its positive
primitive is LC at EVERY centre at least two.
These assertions do not include ordinary LC at centre one, which
does fail: $V=(3,3),u=3$ has first ordinary minor $62-41\phi<0$.
ONE-FACTOR tails are deliberately excluded.

For $u\ge16$, remove a LARGEST interval $[t]$ and call the rest
$L_0$ of length $r_0\ge1$. Then
$u=\deg L_0+t-2\ge\deg L_0+1$, so its OLD quotient is registered
positive global LC with the entire later-primitive property.
Its window gives ordinary LC of $J$ at all centres $n\ge t$,
and primitive LC at all centres $n\ge t+1$.
The new ordinary head $J_n=(LQ_u)_n-c\psi^n$ through rank $u$
is positive, since $L_n\ge1$ there and $|c|<1$.
The relaxed fourfold budget pays every ordinary centre $2..u-1$;
$[z]L=r\le(u+1)/2$. Since $t\le u$ for a nonempty old tail,
this joins the entire ordinary tail after deleting rank zero.
Also $J_0<2$, $J_1\ge u+r-2q>u$, so $qJ_1>J_0$ at $u\ge16$.
The first increasing edge and LC from centre two make $J$ unimodal.

For its primitive put $B=L_0Q_u$, $a=\phi^2c<0$,
$X_n=B_n-a\psi^n$, $b=B_1=u+r_0-q$.
Now $|a|/q\le4$, so the SAME fourfold head pays every required
$X$ minor through $u-1$. Its entries at ranks $1..u$ are positive:
$|a|\le4q$ gives $|a\psi^n|\le4q^2<2$, whereas the read entries
of $B$ are at least $u-q$ through $u-2$ and at least $3/2$ at
the last two ranks; at those last ranks the atom is much smaller
than one. At rank one, $B_1\ge u+1-q>u$.
As before $d_n=X_n$ for $1\le n<t$, $d_t=X_t-1$ and
$d_{t+1}=X_{t+1}-b$.

For $t\le u-1$ set $D=b^2-2B_2=u-r_0-q^2$.
The concave scalar
$F(r)=(u-r-1)(r+1)(u-1)-u(u+r)$ is nonnegative on
$1\le r\le(u-1)/2$: its endpoints are
$u^2-7u+4$ and $(u^3-7u^2+u+1)/4$.
Thus the unperturbed last-junction gap is at least $u$.
Its perturbation is bounded by
\[
 |a|q^{t-1}(b+2q)\le\frac{16}{25}
                         \left(\frac{3u}{2}+\frac16\right)
                         \le\frac{39u}{40}.          \tag{2}
\]
For the first inequality: if $t\ge4$, use $|a|\le4q$ and
$4q^4\le16/25$. If $t=3$ is LARGEST, all intervals are threes,
so $u=2r_0+1\ge16$ gives $r_0\ge3$; hence
$|a|q^2\le16q^7\le16/25$, using $q^5\le1/10$ and
$q^2\le2/5$. The perturbed gap is at least $u/40$.
The final primitive minor is consequently at least $1+u/40$.

The exceptional long junction $t=u$ forces $L_0=[3]$. Exact division
then gives
\[
             d_{u+1}=-u+2q(1-\psi^u)<0.             \tag{3}
\]
Indeed $B_{u+1}=\phi-2q^{2u+1}$ in formal division and
$a\psi^{u+1}=-2q^{2u+1}+2q\psi^u$.
The positive primitive therefore ends by $u$; its last centre is
an endpoint, while centre $u-1$ was already paid by lowering the
next $X_u$ by one. No unproved final robust row at centre $u$
is used. This completes the unbounded induction.
The COMPLETE small range $u\le15$, $h=u+1$, $r\ge2$ has
215 contexts, all independently rebuilt and kernel certified.

\paragraph{The actual boundary unit aligns with the geometric inverse.}
At $u=h-1$ the infinite coefficient tail $C=L\sum_{k=0}^{u-1}
(z+z^2)^k$ satisfies $x_s=C_{u+s}$ for every $s\ge0$.
The missing leading coefficient of $L$ gives EXACTLY
\[
       x_{-1}=C_{u-1}+1.                              \tag{4}
\]
This is the actual derivative at physical rank $W$; for every
$s\ge0$, physical rank $W+1+s$ has positive derivative
$(x_s-E_{s-a})^+$, $a=\alpha-1$, with the true source
$E=(1+qz)J+c$.
The extra actual unit must NOT be discarded: unlike the $u=h$
case, it is exactly what completes the geometric inverse here.

By (1), $x$ is globally LC and all its ratios are below $\phi$
from rank zero. The relaxed unimodal-$b$ cut theorem pays centres
$a+2$ through the first $J$ mode. Source LR, finite support and
$x_M<\phi J_M$ put every pre-first-$A$-mode centre in that range.
The registered first/second source startup payments apply when
$w=u+r\ge8$. The remaining $w<8$ has exactly FOUR tail contexts
and SEVENTEEN paddings $1\le\alpha\le u$. The padding clamp follows
from the positive symmetric-row centre bound
$\operatorname{firstmode}C\le2u-1$; every $\alpha\ge u+1$ starts
after it. All finite startup comparisons are independently rebuilt from
positive suffix-box sums of literal actual coefficients. Thus the result is actual
LC at every centre $W+1\le N<\ell$, whenever this pre-mode band
is nonempty. There is no assertion at the earlier centre $W$.

\paragraph{C uses a different ordinary seed, with the correct positive Cassini.}
The first defect is now $N_0=W+1$. Let $b_*=r_W$, $a_*=r_{W+1}$,
$d_*=r_{W+2}=a_*+b_*$. Its exact conjugate error is
$a_*-\phi b_*=c<0$, so
$\lambda_0=(a_*-\delta)/b_*<\phi$, $\delta\in\{1,2\}$.
Every later increasing current ratio is at most $\lambda_0$ by
the proved actual LC. All cumulative reads through $W$ are
unchanged reference reads and have ratios greater than $\phi$.
The required ordinary seed is at $W$, since
$F_{W-1}=a_*$ and $F_W=a_*+b_*$.
The homogeneous Cassini identity at centre $W+1$ yields
$b_*^2+a_*b_*-a_*^2=A_V$, so its EXACT seed margin is
\[
 b_*(a_*+b_*)-(a_*-\delta)a_*
                 =A_V+\delta a_*>0.                \tag{5}
\]
This does NOT incorrectly invoke a homogeneous slope recurrence at
$W-1$, where the numerator's last coefficient still matters.
The ordinary identity is $F_j=r_{j+2}-B_{j+2}$; it already has
zero correction at $j=W-1$. The separate boundary-read correction
record withdraws an unused aside in the earlier immutable proof.
Ordinary/cumulative integration from this seed covers every later
read. Homogeneous reference C at $N_0=W+1$, improved by the first
defect, pays the first current comparison. This proves ALL
late actual C for $u=h-1$ with $r\ge2$; the one-factor theorem
and registered $u\ge h$ results complete the asserted union.

\paragraph{Further structural observation, not used above.}
For ANY deficit $d=h-u$, the same missing-leading-coefficient
calculation gives model/actual alignment at physical $W$:
$x_{d-2}=C_{h-2}+1$ and $x_s=C_{u+s}$ for $s\ge d-1$.
Thus the actual positive derivative tail from $W$ is always the
geometric quotient cut at offsets $s\ge d-2$, regardless of the
sign of the early quotient. Larger deficits require additional
quotient-head deletion and startup payments and remain OPEN.

\paragraph{Verification and scope.}
The authoritative registration is
\texttt{subdegree\_CI\_entrance\_verification.json}.
\texttt{SubdegreeCI.lean} has eleven new scoped kernel reports;
\texttt{SubdegreeCIChunkCertificate.lean} has three finite reports,
including one axiom-free integer certificate. Ordinary kernel
reduction over 31 bounded chunks checks all 11853 distinct field
margins and all 24 distinct integer startup margins. A separate
capacity-multiset grammar and reverse polynomial division rebuild
all 215 low-degree quotient contexts. Independent positive suffix-
box sums rebuild all 17 small actual hosts and all 31 startup minors,
without using the iterated append algorithm of the primary QA.
The unbounded ULC/window/convolution and unimodal-cut assemblies,
model identification and all-read integration are written proofs,
not an end-to-end Lean theorem for arbitrary compositions.

The one-factor $[6],u=4$ primitive minor $4-3\phi<0$ remains an
explicit exclusion control; its actual C is already known by the
separate one-jump theorem. The first ordinary minor of the actual
multifactored quotient $[3]^2,u=3$ is genuinely negative, so the
weaker inverse/cut argument is necessary. The pre-mode LC assertion
here is ONLY for the new count $u=h-1$. It is not propagated to the
same physical centre at $u=h$, where the prior $(8,12,19)$ actual
LC counter remains valid. The union is a theorem of C comparisons.
Counts $0\le u<h-1$, earlier reference C, general C/A2 and the
ORIGINAL general unamplified one-descent UI remain OPEN. No original
least-counterexample numerical bound is upgraded.
